
\begin{appendices}

\section{支撑材料列表}

% 以下是支撑材料列表的示例
\begin{itemize}\ttfamily
  \item{database.xlsx}
  \item{dataset.csv}
  \item{algorithm\_a.cpp}
  \item{algorithm\_b.m}
  \item{algorithm\_c.py}
\end{itemize}

% windows系统下的电脑有个导出技巧：
%% 1. 来到当前文件夹页面
%% 2. 全选中所有文件
%% 3. 点击文件夹页面中“主页”选项卡的“复制路径”
%% 4. 在空白处粘贴
%% 5. 利用全局替换的方式将根目录统一替换成\item{，然后再单独加上右花括号
%% 6. 做最后编辑（尤其是下划线要加上反斜杠，如上述示例所示），即可得到itemize环境的内容


\section{相关定理的证明}

\zhlipsum[2]


\section{程序代码}

% 方式一：直接输入的形式插入代码

\begin{lstlisting}[language=C++,title={\raggedright 基础程序}]
  int main(){
      int i;
      printf("hello latex!\n");
      return 0;
  }
\end{lstlisting}


% 方式二：以文件形式插入代码

\lstinputlisting[ language=C++, title={\raggedright 计算$\boldsymbol{n}$的阶乘（C++）：} ]{codes/funfactorial.cpp} %C++

\lstinputlisting[ language=Java, title={\raggedright 计算$\boldsymbol{n}$的阶乘（Java）：} ]{codes/funfactorial.java} %Java

\lstinputlisting[ language=Python, title={\raggedright 计算$\boldsymbol{n}$的阶乘（Python）：} ]{codes/funfactorial.py} %Python

\lstinputlisting[style=Matlab-editor, title={\raggedright 计算$\boldsymbol{n}$的阶乘（MATLAB）：}]{codes/funfactorial.m} %MATLAB（如果不想要这种风格，则把该行命令的可选参数style=Matlab-editor改为language=Matlab

\end{appendices}
